\subsection{考虑均值的表示}

中心化结果提示，应研究主码符号何时应编码相对均值的偏差（\cref{sec:8-centring}）。首先需要对中心化和未中心化数据重复冷建，以区分其效果与训练波动。列表相关的均值也是候选方向，但要求查询评分与残差重建保持一致。测试应比较短名单覆盖率、可达到的最佳精排召回率、所选预算和总吞吐，防止较便宜的短名单掩盖质量上限下降。

\subsection{校准风险与复用成本}

具有统计控制的预算选择，能否省下比获取预算更大的查询时间？在 \cref{eq:6-rank-curve} 的嵌套条件下，参考近邻的主码排名可代替反复试搜。早期四配置离线回放发现，序贯抽样增加了搜索次数；用已测校准时间拟合的模型估算，即使样本更少，成本仍然更高。这不是完整 GPU 控制器的测量。替代方案应测试批量试验或排名复用，并以直接校准加复用总时间，对比质量匹配的固定预算。风险分析必须覆盖抽样和预算选择程序，发生漂移时应更新校准；相对参考的控制仍需与标签建立联系，才能给出绝对召回保证。

\subsection{GPU 与主机内存间的残差存储}

BGE 和 Stella 的残差占进程总常驻显存的 80.7--82.9\%（\cref{sec:7-7}），但只有短名单中的行需要读取。将路由和主码留在 GPU，按块从主机内存取回残差，有望增加容量。难点在于把分散的短名单位置转成高效块传输。在评分与选择固定时，改变设备容量上限，测量额外字节、缓存命中、传输重叠、吞吐和尾延迟。由此才能确定容量与延迟的取舍，而非假定卸载会提高速度。

\subsection{解释小 batch 的边界}

在召回率 0.97 处，分组改善了 JHQ 的小 batch 执行，却没有消除与 CAGRA fp32 的差距（\cref{sec:7-smallbatch}）。候选工作量、选择开销和接口成本都可能参与其中。匹配的研究应统计两种索引评分的候选数，并区分内核时间、传输和主机调用开销。跨 batch 比较这些成本，才能检验下一步应减少已探测候选、降低选择成本，还是使用不同索引。
